%! Author = lili
%! Date = 6/14/26


%========================= GRUNDLAGEN =========================


\section{0 -- Grundbegriffe \& Modellorganismus}
\emph{A. thaliana:} Genom ca. 125 Mb, 5 Chromosomen, ca. 27000 Gene; klein, schnell, selbstbestäubend (Zwitterblüten), leicht transformierbar $\Rightarrow$ ideales Modell.
\begin{itemize}
    \item \textbf{Vorwärtsgenetik:} Phänotyp $\to$ Gen (erst Mutante, dann Gen; z.B. EMS, V5).
    \item \textbf{Rückwärtsgenetik:} Gen $\to$ Phänotyp (gezielte Mutation; z.B. T-DNA/CRISPR, V1).
    \item \textbf{homozygot/heterozygot:} beide/zwei versch. Allele. Meiste Mutationen \textbf{rezessiv} $\Rightarrow$ Phänotyp nur homozygot.
    \item \textbf{CRISPR/Cas9:} gRNA lenkt Cas9 zur Zielsequenz $\to$ Doppelstrangbruch $\to$ fehlerhafte Reparatur (InDels) $\to$ Gen zerstört. Wird per T-DNA eingebracht.
\end{itemize}


%========================= V1 =========================


\section{V1 -- Genotyping mittels PCR (T-DNA/CRISPR)}
\textbf{PCR-Prinzip (3 Schritte/Zyklus):} Denaturierung (94\,°C) $\to$ Annealing (56--60\,°C, Primer binden) $\to$ Elongation (72\,°C, Taq). n Zyklen $\to$ ca. 2\,hoch\,n Kopien.
\begin{itemize}
    \item \textbf{Kontrollen:} Positivkontrolle (unabh. Gen, hier S4/S7) $\Rightarrow$ DNA amplifizierbar. Negativkontrolle (H\textsubscript{2}O) $\Rightarrow$ deckt Kontamination auf.
    \item \textbf{Gelelektrophorese:} DNA negativ $\to$ wandert zur \textbf{Anode}; Agarose = Sieb, klein läuft weiter. EtBr interkaliert, UV-sichtbar (mutagen $\to$ Handschuhe!).
    \item \textbf{NcoI/CAPS-Prinzip:} WT-Allel hat NcoI-Site, Mutante nicht $\to$ Verdau: WT geschnitten, Mutante ungeschnitten.
\end{itemize}
\textbf{Bandenmuster T-DNA (3-Primer LP/RP/BP):}
\begin{itemize}
    \item WT: nur WT-Bande (LP+RP).
    \item hom. Insertion: nur Insertionsbande (BP+RP).
    \item heterozygot: \textbf{beide} Banden.
\end{itemize}


\kfrage{bandenmuster}
\frage{genotypphaenotyp}
\frage{vorwaertsrueckwaerts}
\frage{crispr}
\frage{arabidopsismodell}
\frage{kontrollenpcr}
\frage{gelelektrophorese}
\frage{ncoicaps}

%========================= V2 =========================


\section{V2 -- GUS-Reportergen}
\textbf{Reporter:} \textbf{uidA} (E.\,coli) $\to$ \textbf{$\beta$-Glucuronidase (GUS)}. Substrat \textbf{X-Gluc} $\to$ blaues Indigo-Präzipitat am Ort der Aktivität. Kein Hintergrund in Pflanzen, stabil, fixierbar.
\begin{itemize}
    \item \textbf{Schritte:} Fixierung (Formaldehyd) $\to$ Waschen $\to$ X-Gluc + Vakuuminfiltration, ü.N. 37\,°C $\to$ Entfärben mit EtOH (Chlorophyll raus) $\to$ Mikroskopie.
    \item \textbf{transkriptionelle Fusion:} Promotor + Reporter $\to$ misst \emph{Promotoraktivität} (wo transkribiert).
    \item \textbf{translationale Fusion:} Reporter im Leseraster am Gen $\to$ misst zusätzlich \emph{Proteinlokalisation}.
    \item \textbf{Konstrukte:} pro35S (CaMV, konstitutiv, \textbf{Positivkontrolle}), proMYB12 (\textbf{Wurzel}), proMYB111 (\textbf{Kotyledonen}), Col-0 = keine Färbung.
\end{itemize}


\kfrage{gusschritte}
\kfrage{gusreporter}
\kfrage{guspromotor}
\kfrage{gusfärbung}
\kfrage{fusionunterschiede}

%========================= V3 =========================


\section{V3 -- Genexpression via RT-PCR}
\textbf{Schritte:} RNA-Extraktion $\to$ OD\textsubscript{260} messen $\to$ \textbf{DNase-Verdau} $\to$ Reverse Transkription (cDNA) $\to$ PCR $\to$ Gel.
\begin{itemize}
    \item \textbf{Warum cDNA?} Taq braucht DNA-Matrize; RNA muss erst umgeschrieben werden.
    \item \textbf{DNase-Verdau:} entfernt genomische DNA, sonst falsch-positive PCR.
    \item \textbf{Aktin2-Kontrolle:} Housekeeping-Gen $\Rightarrow$ cDNA ok $\Rightarrow$ Fehlen von MYB ist echt.
    \item \textbf{OD:} 1 OD\textsubscript{260} = 40\,µg/ml RNA. Reinheit A260/A280 $\approx$ 2,0 (RNA), 1,8 (DNA).
    \item \textbf{Häufigkeit:} rRNA $\gg$ tRNA $\gg$ mRNA (nur 1--4\,\%). Bleach-Gel: 28S/18S sichtbar.
\end{itemize}
\textbf{RT-Primer (3 Arten):} oligo-dT (nur mRNA, ab 3'), Random-Hexamer (gesamte RNA, kurze cDNAs), genspezifisch (nur Zieltranskript).


\kfrage{rtpcr}
\kfrage{rtpcrprimer}
\frage{rnaalspcr}
\frage{aktinkontrolle}
\frage{od260}
\kfrage{rnahäufigkeit}
\kfrage{genexpressiontranskript}
\kfrage{genexpressionsanalyse}
\kfrage{transkriptionbistranslation}
\kfrage{zwischentranskriptionundtranslation}


%========================= mRNA-Reifung =========================


\section{Transkription $\to$ Translation (mRNA-Reifung)}
\textbf{5 RNA-Stufen / 4 Schritte:}
\begin{enumerate}
    \item pre-mRNA (hnRNA)\quad$\downarrow$ \textbf{Capping} (5'-m7G)
    \item gecappte RNA\quad$\downarrow$ \textbf{Spleißen} (Introns raus, Exons ligiert)
    \item gespleißte RNA\quad$\downarrow$ \textbf{Polyadenylierung} (Poly-A)
    \item polyadenylierte RNA\quad$\downarrow$ \textbf{Kernexport}
    \item reife mRNA $\to$ Translation
\end{enumerate}


%========================= V4 =========================


\section{V4 -- Genetische Kartierung}
\textbf{Genetische Karte:} relative Lage/Abstand von Markern. Varianten: \textbf{Kopplungskarte} (cM, aus Rekombination) vs. \textbf{physikalische Karte} (bp, aus Sequenz). 1\,cM = 1\,\% RF.
\begin{itemize}
    \item \textbf{DFLP/InDel-Marker:} versch. \emph{Länge} (z.B. TT5: Col 408\,bp, Ler 456\,bp).
    \item \textbf{CAPS-Marker (PCR-RFLP):} versch. \emph{Schnittstelle} (TT3: AflII nur in Col $\to$ 630+450\,bp; Ler 1,1\,kb).
    \item \textbf{heterozygot} $\to$ zeigt \textbf{beide/alle} Banden.
\end{itemize}
\textbf{Rekombination zählen (F2 Ler$\times$Col), Genotypen L/C/H:}
\begin{itemize}
    \item H\,$\leftrightarrow$\,L oder H\,$\leftrightarrow$\,C = \textbf{1 Crossover}
    \item L\,$\leftrightarrow$\,C = \textbf{2 Crossover}
\end{itemize}
$\text{RF} = \dfrac{\text{rekombinante Chromosomen}}{\text{Pflanzen}\times 2}$.\;
Kleine RF $\to$ enge Kopplung; RF $\approx$ 50\,\% $\to$ ungekoppelt.


\kfrage{genetKarten}
\kfrage{rekomnbinationsfrequenzen}
\kfrage{kreuzungauswertung}
\kfrage{rekombinatenmarker}
\kfrage{markertypen}
\kfrage{deletionenfinden}
\kfrage{sequenzenunterscheiden}
\frage{segregation}

%========================= V5 =========================


\section{V5 -- Metabolitanalyse (HPTLC / Flavonoide)}
\textbf{HPTLC:} hochauflösende Dünnschichtchromatographie. Kieselgel (stationär, polar), Laufmittel steigt kapillar auf, Trennung nach Polarität; Detektion mit DPBA + PEG unter UV. \emph{HPTLC vs. TLC:} feineres Gel $\to$ bessere Auflösung/Empfindlichkeit/Reproduzierbarkeit.
\begin{itemize}
    \item \textbf{Vorteile:} schnell, günstig, viele Proben parallel. \textbf{Nachteile:} geringe Auflösung/Sensitivität, kaum quantitativ.
    \item \textbf{Auswertung:} Rf-Wert + Fluoreszenzfarbe vs. Standards. Kaempferol $\to$ grün/gelb-grün, Quercetin $\to$ orange. Fehlende Bande = Weg blockiert.
\end{itemize}


\section{Flavonoid-Stoffwechselweg (TT-Gene)}
\centering
\begin{tikzpicture}[
    node distance=3.5mm,
    every node/.style={font=\scriptsize},
    enz/.style={midway,right,font=\scriptsize\bfseries},
    ->,>={Stealth[length=2mm]}]
    \node (a) {p-Cumaroyl-CoA};
    \node (b) [below=of a] {Naringenin-Chalkon};
    \node (c) [below=of b] {Naringenin};
    \node (d) [below=of c] {Dihydroflavonol};
    \node (e) [below=of d] {\textbf{Flavonole} (Kaempferol/Quercetin)};
    \node (f) [below=of e] {Anthocyane / Proanthocyanidine};
    \draw (a) -- node[enz]{CHS (TT4)} (b);
    \draw (b) -- node[enz]{CHI (TT5)} (c);
    \draw (c) -- node[enz]{F3H (TT6)} (d);
    \draw (d) -- node[enz]{FLS} (e);
    \draw (d.west) to[bend right=55] node[left,font=\scriptsize\bfseries]{DFR (TT3)} (f.west);
\end{tikzpicture}
\vspace{2pt}

\raggedright\textbf{Regulation:} MYB12/MYB111/MYB11 aktivieren frühe Gene (CHS, CHI, F3H, FLS $\to$ Flavonole). MBW-Komplex TT2(MYB)/TT8(bHLH)/TTG1(WD40) reguliert späte Gene.


\kfrage{HPTCL}
\kfrage{durchführungPTCL}
\kfrage{HPTCLvsPTCL}
\kfrage{ptclvorteile}
\kfrage{HPTLCauswertung}
\kfrage{stoffwechselwegmutante}


%========================= Verknuepfung =========================


\section{Roter Faden: MYB12 vs. MYB111}
Konsistent über alle Versuche -- flavonol-spezifische R2R3-MYB-Aktivatoren:
\begin{itemize}
    \item \textbf{MYB12} aktiv in der \textbf{Wurzel}.
    \item \textbf{MYB111} aktiv in den \textbf{Kotyledonen}.
\end{itemize}
$\Rightarrow$ GUS (blau in Wurzel bzw. Kotyledonen) = RT-PCR (cDNA in Wurzel bzw. Kotyledonen) = HPTLC (gewebespez. Flavonolmuster). \emph{Ideale Klausur-Synthesefrage.}


\frage{myb12myb111}

%========================= Bioinformatik =========================


\section{Bioinformatik}
\begin{itemize}
    \item \textbf{Assemblierung:} überlappende Reads $\to$ Contigs $\to$ Scaffolds. de-novo vs. referenzbasiert.
    \item \textbf{N50:} Länge, bei der 50\,\% des Assemblies in Contigs $\geq$ dieser Länge liegen. \textbf{Hoch = besser} (kontiguierlicher), aber sagt nichts über Korrektheit.
    \item \textbf{Repeats:} Reads kürzer als Repeat $\to$ nicht eindeutig $\to$ Kollaps, Missassembly, Lücken. Lösung: Long Reads / Paired-End.
    \item \textbf{Gene finding:} ab initio (ORFs, Codon-Usage, GC, Spleiß-Stellen, HMM) + evidenzbasiert (BLAST, RNA-Seq, Konservierung).
    \item \textbf{Kodierend?} lange ORFs, Codon-Bias, Homologie, RNA-Seq-Nachweis, Ka/Ks\,$<$\,1.
    \item \textbf{lokal vs. global:} Needleman-Wunsch (global, Ende-zu-Ende) vs. Smith-Waterman (lokal, beste Teilbereiche).
    \item \textbf{BLAST vs. SW:} SW = optimal aber langsam O(n·m); BLAST = heuristisch (Seeds/k-mere, dann Extension) $\to$ viel schneller, fast optimal.
\end{itemize}


\kfrage{assemblierung}
\kfrage{bioinf}
\kfrage{repeatprobleme}
\kfrage{genefinding}
\kfrage{codingfinden}
\kfrage{blastsmithwaterman}
\kfrage{alignmentunterschiede}
\kfrage{multiplesalignment}

%========================= FORMELN =========================
    \begin{itemize}
        \item Verdünnung: $c_1 V_1 = c_2 V_2$
        \item Stoffmenge: $n = \dfrac{m}{M}$ \; ; \; $m = n\cdot M$ \; ; \; $c = \dfrac{n}{V}$
        \item Prozent: (w/v) = g/100\,ml \; ; \; (v/v) = ml/100\,ml
        \item RNA: $c = \text{OD}_{260}\cdot 40\,\tfrac{\mu g}{ml}\cdot \text{Verd.}$
    \end{itemize}
    \textbf{Fallen:} mM $\leftrightarrow$ M (1\,M = 1000\,mM) umrechnen; gefragte Größe genau lesen; überflüssige/irreführende Angaben; Einheiten mitführen!

\kfrage{konzentrationsberechnung}
\kfrage{rechenaufgaben}
\kfrage{rechenaufgabenII}


